


Lorsque \(X\) est propre sur \(k\), on sait [4] que \(\text{\uwave{Pic}}^0_{X/k}\) est représentable par un schéma en groupes localement de type fini \(\underline{\Pic}_{X/k}\) sur \(k\). La composante neutre \(\underline{\Pic}^0_{X/k}\) de \(\underline{\Pic}_{X/k}\) est un groupe algébrique commutatif connexe sur \(k\).

\textbf{Définition 2.4.1.} \uline{Soit} \(X\) \uline{un schéma propre sur un corps} \(k\), \(\underline{L}\) \uline{un} \(\underline{O}_X\)\uline{-module inversible} (resp. \(x\) \uline{un élément de} \(\Pic(X)\)\uline{). On dit que} \(\underline{L}\)(resp. \(x\)\uline{) est} \(\tau\)\uline{-équivalent à zéro, s'il existe un entier} \(n\ne0\) \uline{tel que} \(\underline{L}^{\otimes n}\)(resp. \(nx\)\uline{) est algébriquement équivalent à zéro. Les éléments de} \(\Pic(X)\) \uline{algébriquement équivalents} (resp. \(\tau\)\uline{-équivalents) à zéro forment un sous-groupe} \(\Pic^0(X)\) (resp. \(\Pic^\tau(X)\)\uline{) de} \(\Pic(X)\).

\textbf{Lemme 2.4.3.} \uline{Pour tout schéma propre} \(X\) \uline{sur un corps algébriquement clos} \(k\) \uline{le groupe} \(\Pic^0(X)\) \uline{est divisible par tout entier} \(n\) \uline{premier à la caractéristique de} \(k\).

\emph{Preuve} : Comme \(k\) est algébriquement clos, on a \(\Pic^0(X)\cong \underline{\Pic}^0_{X/k}(k)\). Puisque \(\underline{\Pic}^0_{X/k}\) est connexe et de type fini, l'endomorphisme puissance \(n\)-ème est surjectif lorsque \(n\) est premier à la caractéristique de \(k\) (SGA 3 XV 1.3 et VI 1.3.2). Mais \(k\) est algébriquement clos, donc l'endomorphisme correspondant sur le groupe \(\underline{\Pic}^0_{X/k}(k)\) est surjectif, C.Q.F.D.

De 2.4.3 on conclut que dans la suite exacte
\[
0\longrightarrow\Pic^0(X)\longrightarrow\Pic^\tau(X)\longrightarrow T\longrightarrow0,
\]
\(\Pic^0(X)\) est divisible par un nombre premier \(p\) et \(T\) est un groupe de torsion, de sorte qu'on peut appliquer 2.1.3 à la restriction de l'application polynôme (2.3.1) à \(\Pic^\tau(X)\). On trouve ainsi :

